\section*{Problem \#302}

\subsection*{1) ROUND-2 OBJECTIVE}
\textbf{Path C (obstruction/correction).}
Round~1 already showed that the specific conjecture
\[
  f(N)=(1/2+o(1))N
\]
is false by proving the explicit lower bound \(f(N)\ge (5/8+o(1))N\).
The most promising round-2 task is therefore not to re-open the false \(1/2\)-claim, but to sharpen the corrected picture.
In this round I prove a new \emph{family} of solution-free sets, containing the parity construction as a special case, and I show that within this whole family the density \(5/8\) is already optimal.

\subsection*{2) Round-1 FOUNDATION USED}
I use the following Round~1 results without re-proving them.
\begin{enumerate}
\item If \(1\le a,b,c\le N\) are pairwise distinct and \(1/a=1/b+1/c\), then \(a<N/2\).
\item The identity \(bc=a(b+c)\) is equivalent to \(1/a=1/b+1/c\).
\item The Round~1 construction
\[
  [\lceil N/2\rceil,N]\ \cup\ \{n\le N/4:\ n\text{ odd}\}
\]
is solution-free and has size \((5/8+o(1))N\).
\end{enumerate}

\subsection*{3) NEW INSIGHT / TOOL (ROUND-2)}
The new tool is a \emph{single-residue-class modular obstruction}.
Instead of taking the low block to be the odd numbers, I take it to be one fixed invertible residue class modulo an arbitrary modulus \(m\ge 2\).
The parity construction is exactly the case \((m,r)=(2,1)\).
This gives a larger theorem and, after optimizing the interval parameter, proves that \(5/8\) is the best density obtainable from any construction of this one-residue/one-interval type.

\subsection*{4) ATTACK PLAN (ROUND-2)}
The gap left after Round~1 was that the \((5/8)\)-construction looked somewhat ad hoc.
To improve on that, I will:
\begin{enumerate}
\item prove a general theorem for sets of the form
\[
  \{n\le \beta N/2:\ n\equiv r\pmod m\}\cup [\beta N,N],
\]
with \(1/2\le \beta\le 1\) and \((r,m)=1\);
\item optimize the resulting density over \(\beta\) and \(m\);
\item check carefully that the modular argument still works when \(m\) is composite and one of \(b,c\) is not coprime to \(m\).
\end{enumerate}

\subsection*{5) WORK (ROUND-2)}
\textbf{Theorem.}
Fix integers \(m\ge 2\) and \(r\) with \((r,m)=1\), and fix a real number \(\beta\) with \(1/2\le \beta\le 1\).
Define
\[
  L_N(m,r,\beta):=\{n\in\{1,\dots,N\}: n\le \lfloor \beta N/2\rfloor,\ n\equiv r\pmod m\},
\]
\[
  U_N(\beta):=\{\lceil \beta N\rceil,\lceil \beta N\rceil+1,\dots,N\},
\]
and
\[
  A_N(m,r,\beta):=L_N(m,r,\beta)\cup U_N(\beta).
\]
Then \(A_N(m,r,\beta)\) contains no pairwise distinct \(a,b,c\) with
\[
  \frac1a=\frac1b+\frac1c.
\]
Consequently,
\[
  f(N)\ge \left(1-\beta+\frac{\beta}{2m}\right)N+O(1).
\]
Within this family the density is maximized at \(\beta=1/2\), where it becomes
\[
  \frac12+\frac{1}{4m}.
\]
Hence the best member of this entire family is the parity construction \((m,r,\beta)=(2,1,1/2)\), which recovers density \(5/8\).

\medskip
\textit{Proof.}
Assume for contradiction that distinct \(a,b,c\in A_N(m,r,\beta)\) satisfy \(1/a=1/b+1/c\).
By the Round~1 lemma, the smallest denominator in any such triple satisfies \(a<N/2\).
Since \(\beta\ge 1/2\), every element of \(U_N(\beta)\) is at least \(N/2\), so necessarily
\[
  a\in L_N(m,r,\beta).
\]
In particular,
\[
  a\le \beta N/2,
  \qquad a\equiv r\pmod m,
  \qquad (a,m)=1.
\]

I now show that neither \(b\) nor \(c\) can lie in \(L_N(m,r,\beta)\).
Suppose first that \(b\in L_N(m,r,\beta)\).
Then also \(b\equiv r\pmod m\) and \((b,m)=1\).
There are two cases.

\smallskip
\emph{Case 1: \((c,m)=1\).}
Then \(a,b,c\) are all invertible modulo \(m\), and from
\[
  bc=a(b+c)
\]
we may divide by \(abc\) modulo \(m\) to get
\[
  a^{-1}\equiv b^{-1}+c^{-1}\pmod m.
\]
But \(a\equiv b\equiv r\pmod m\), so \(a^{-1}\equiv b^{-1}\equiv r^{-1}\pmod m\), hence
\[
  r^{-1}\equiv r^{-1}+c^{-1}\pmod m,
\]
which gives \(c^{-1}\equiv 0\pmod m\), impossible.

\smallskip
\emph{Case 2: \((c,m)>1\).}
Choose a prime \(q\mid (c,m)\).
Since \(a\equiv b\equiv r\pmod m\) and \((r,m)=1\), we have \(q\nmid a\) and \(q\nmid b\), whereas \(q\mid c\).
Reducing
\[
  bc=a(b+c)
\]
modulo \(q\), the left-hand side is \(0\), while the right-hand side is congruent to \(ab\not\equiv 0\pmod q\).
Contradiction.

Thus \(b\notin L_N(m,r,\beta)\).
By symmetry \(c\notin L_N(m,r,\beta)\) as well.
Therefore both \(b\) and \(c\) must lie in \(U_N(\beta)\), so
\[
  b,c\ge \beta N.
\]
Assume \(b\le c\).
Then, because the denominators are distinct,
\[
  \frac1a=\frac1b+\frac1c<\frac{2}{b}\le \frac{2}{\beta N},
\]
so
\[
  a>\beta N/2.
\]
But we already know \(a\le \beta N/2\), contradiction.
This proves that \(A_N(m,r,\beta)\) is solution-free.

For the size estimate, we have
\[
  |U_N(\beta)|=(1-\beta)N+O(1),
\qquad
  |L_N(m,r,\beta)|=\frac{\beta}{2m}N+O(1),
\]
so
\[
  |A_N(m,r,\beta)|=\left(1-\beta+\frac{\beta}{2m}\right)N+O(1).
\]
Since \(m\ge 2\), the coefficient of \(\beta\) here is
\[
  -1+\frac{1}{2m}<0,
\]
so the density is maximized when \(\beta\) is as small as allowed, namely \(\beta=1/2\).
This gives the best density inside the family as
\[
  \frac12+\frac{1}{4m}.
\]
Finally, this is maximized over \(m\ge 2\) at \(m=2\), giving \(5/8\).
\hfill \(\square\)

\subsection*{6) ADVERSARIAL VERIFICATION}
\begin{enumerate}
\item \textbf{Could the modular step fail for composite \(m\)?}
No. If \((c,m)=1\), inverses modulo \(m\) exist and the proof is direct. If \((c,m)>1\), one passes to a prime divisor \(q\mid (c,m)\) and obtains a contradiction modulo \(q\).
\item \textbf{Could \(b\) and \(c\) both be in the upper block without contradiction when \(b=\beta N\)?}
No. Distinctness gives
\[
  \frac1b+\frac1c<\frac2b\le \frac{2}{\beta N},
\]
so the resulting inequality for \(a\) is strict.
\item \textbf{Did I accidentally use a larger family than I actually proved?}
No. The theorem covers exactly one residue class modulo \(m\) in the low interval, not an arbitrary set of residues. I am not claiming optimality beyond this family.
\item \textbf{Does the density computation double-count an overlap between low and upper blocks?}
No. Since \(\beta/2<\beta\), the intervals \([1,\beta N/2]\) and \([\beta N,N]\) are disjoint for every \(\beta>0\).
\end{enumerate}

\subsection*{7) FINAL (EXACTLY ONE)}
\textbf{UNRESOLVED (BUT STRICTLY ADVANCED).}
Round~2 proves a genuine new theorem: the Round~1 parity construction is the optimal member of the entire one-residue/one-interval family
\[
  \{n\le \beta N/2:\ n\equiv r\pmod m\}\cup [\beta N,N].
\]
This does not determine the true asymptotic of \(f(N)\), but it does explain rigorously why the Round~1 \((5/8)\)-construction is the best possible within a broad modular template.

\subsection*{8) COMPLETION ESTIMATE}
\textbf{COMPLETION: 40\%}

\subsection*{9) REFERENCES}
\begin{itemize}
\item Round~1 output in this conversation (used as fixed foundation).
\item T. F. Bloom, \emph{Erd\H{o}s Problem \#302}.
\end{itemize}


